\section{Демоническое программирование}

\subsection{Домашнее задание}

\begin{enumerate}

  \item
    Дан набор нечестных монеток с вероятностью выпадения орла
    $p_1, p_2, \dots, p_n$. Требуется посчитать вероятность выпадения
    ровно $k$ орлов за $\O(n \cdot k)$. Операции над числами считать
    выполнимыми за $\O(1)$.

  \item
    Дана строка $s$ длины $n$. Для каждой пары $(i, j)$ найти длину максимального
    общего префикса $i$-го и $j$-го суффиксов строки $s$. $\O(n^2)$.

  \item
    Дана строка из латинских букв длины $n$, нужно ее за $\O(n^3)$ запаковать в максимально короткую, используя правило
    $n(S) = \underbrace{SS{\dots}S}_{n}$. \\ Например \texttt{UNIXYESYESYESUNIXYESYESYES} $\rightarrow$ \texttt{2(UNIX3(YES))}.

  \item
    \begin{enumerate}
      \item Даны пара чисел $k, v$ и массив $a$ размера $W$. Постройте за $\O(W)$ массив $b$ такого же размера, где $b_i = \max\limits_{0\leq j \leq \min(k, i)} \Bigl( a_{i - j} + v\cdot j \Bigr)$.
      \item Даны $n$ предметов. У каждого есть цена $v_i$ и вес $w_i$. Каждый предмет можно взять от $0$ до $k_i > 0$ раз. Найдите за $\O(nW)$ максимальную суммарную цену, которую можно набрать таким образом, чтобы суммарный вес не превышал $W$.
    \end{enumerate}

  \item
 	Пусть есть $n$ подарков разной натуральной стоимости и три поросёнка. Нужно
    раздать подарки как можно честнее (так, чтобы минимизировать
    разность суммарной стоимости подарков самого везучего поросёнка и
    самого невезучего). Придумайте алгоритм решения данной задачи за
    $\O(nW^2)$, где $W$ --- суммарная стоимость подарков.

\subsection*{Дополнительные задачи}

  \item
    У профессора есть $k$ яиц и $n$ этажное здание.
    Он хочет узнать максимальное $x$: если яйцо бросить с $x$-го этажа, оно не разобьётся.
    Не разбившиеся яйца можно переиспользовать. Минимизировать число бросков в худшем случае.
    \begin{enumerate}
      \item $\O(n^2 k)$.
      \item $\O(n k \log{n})$.
      \item $\O(n k)$.
      \item $\O(n \log{n})$.
    \end{enumerate}

\end{enumerate}

\newpage

\subsection{Решения}

\begin{enumerate}

    \item

    Обозначим \( dp[i][j] \) как вероятность того, что при первых \( i \) монетках выпадет
    ровно \( j \) орлов. Для каждой монеты \( i \) вероятность того, что она выпадает орлом,
    обозначим как \( p_i \). Следовательно, вероятность того, что она выпадает решкой, будет равна \( 1 - p_i \).

    Для вычисления \( dp[i][j] \) мы можем использовать следующее рекуррентное соотношение:
    \[
        dp[i][j] = p_i \cdot dp[i - 1][j - 1] + (1 - p_i) \cdot dp[i - 1][j]
    \]

    где:
    \begin{itemize}

        \item
            \( p_i \cdot dp[i - 1][j - 1] \) — вероятность того, что \( i \)-я монета выпадет орлом,
            при этом на первых \( i-1 \) монетках уже выпало \( j-1 \) орлов.

        \item
            \( (1 - p_i) \cdot dp[i - 1][j] \) — вероятность того, что \( i \)-я монета выпадет решкой,
            при этом количество орлов остается неизменным.

    \end{itemize}

    Для корректной работы алгоритма необходимо задать начальные условия:
    \begin{itemize}
        \item \( dp[0][0] = 1 \): вероятность получения 0 орлов при 0 бросках равна 1.
        \item \( dp[0][j] = 0 \) для всех \( j > 0 \): вероятность получения более 0 орлов при 0 бросках равна 0, так как отсутствуют броски.
    \end{itemize}

    Искомый ответ будет равен \( dp[n][k] \).

    \item

    Пусть $dp[i][j]$ -- длина наибольшего общего префикса $s[i:]$ и $s[j:]$ суффиксов.
    Пытаемся продлить прошлый общий префикс, если значения в $i$-й и $j$-й позициях совпадают.

    Тогда
    \[
        dp[i][j] =
        \begin{cases}
            dp[i + 1][j + 1] + 1, & \text{если } s[i] = s[j], \\
            0, & \text{иначе}.
        \end{cases}
    \]

    % Дана строка из латинских букв длины $n$, нужно ее за $\O(n^3)$ запаковать в
    % максимально короткую, используя правило $n(S) = \underbrace{SS{\dots}S}_{n}$.
    %\\ Например \texttt{UNIXYESYESYESUNIXYESYESYES} $\rightarrow$ \texttt{2(UNIX3(YES))}.

    \item \dots

    \item \dots

    \item \dots

    \item \dots

\end{enumerate}
